% GENERATED FILE. Edit canonical lessons, then run npm run generate:books.
% canonical-source-hash: fnv1a64:221fc49a9e9e5a27
% canonical-lessons: GE-C146-erklaeren, GE-C146-fehlen, GE-C146-feiern, GE-C146-fuellen, GE-C146-gefallen

\chapter{To explain, To be missing, To celebrate, To fill, To please}
\label{ch:ge-to-explain-to-be-missing-to-celebrate-to-fill-to-please}
\hlchaptermodality{146}

\begin{chapteropening}
\textbf{By the end of this chapter:} \emph{I can say to explain, to be missing, to celebrate, to fill and to please.}

Complete the last lesson of chapter 146: I can say to explain, to be missing, to celebrate, to fill and to please.
\end{chapteropening}

\section[erklären]{erklären — to explain}

\label{lesson:GE-C146-erklaeren}



\begin{quote}
Before the new one: say the German for to recommend, then the German for to decide.
\end{quote}

\subsection*{You'll want to know: erklären}
\textbf{erklären} — \textquotedblleft{}to explain\textquotedblright{}.

\begin{grammarlens}[title={What you've built}]
One more word for this chapter.
\end{grammarlens}

\subsection*{Your turn}
\begin{itemize}
\raggedright
  \item \emph{Say it:} \emph{erklären}
  \item \emph{Say it:} \emph{erklären}, once more
  \item \emph{Say it:} say \emph{erklären}
  \item \emph{Recall:} say the German for to recommend, then the German for to decide, then say \emph{erklären} again
\end{itemize}

\begin{tcolorbox}[breakable,colback=teal!4,colframe=teal!35!black,title={Before you move on}]
What does erklären mean? (\textquotedblleft{}To explain\textquotedblright{}.) Say it once more.
\end{tcolorbox}
\section[fehlen]{fehlen — to be missing}

\label{lesson:GE-C146-fehlen}



\begin{quote}
Before the new one: say the German for to decide, then the German for to explain.
\end{quote}

\subsection*{You'll want to know: fehlen}
\textbf{fehlen} — \textquotedblleft{}to be missing\textquotedblright{}.

\begin{grammarlens}[title={What you've built}]
One more word for this chapter.
\end{grammarlens}

\subsection*{Your turn}
\begin{itemize}
\raggedright
  \item \emph{Say it:} \emph{fehlen}
  \item \emph{Say it:} \emph{fehlen}, once more
  \item \emph{Say it:} say \emph{fehlen}
  \item \emph{Recall:} say the German for to decide, then the German for to explain, then say \emph{fehlen} again
\end{itemize}

\begin{tcolorbox}[breakable,colback=teal!4,colframe=teal!35!black,title={Before you move on}]
What does fehlen mean? (\textquotedblleft{}To be missing\textquotedblright{}.) Say it once more.
\end{tcolorbox}
\section[feiern]{feiern — to celebrate}

\label{lesson:GE-C146-feiern}



\begin{quote}
Before the new one: say the German for to explain, then the German for to be missing.
\end{quote}

\subsection*{You'll want to know: feiern}
\textbf{feiern} — \textquotedblleft{}to celebrate\textquotedblright{}.

\begin{grammarlens}[title={What you've built}]
One more word for this chapter.
\end{grammarlens}

\subsection*{Your turn}
\begin{itemize}
\raggedright
  \item \emph{Say it:} \emph{feiern}
  \item \emph{Say it:} \emph{feiern}, once more
  \item \emph{Say it:} say \emph{feiern}
  \item \emph{Recall:} say the German for to explain, then the German for to be missing, then say \emph{feiern} again
\end{itemize}

\begin{tcolorbox}[breakable,colback=teal!4,colframe=teal!35!black,title={Before you move on}]
What does feiern mean? (\textquotedblleft{}To celebrate\textquotedblright{}.) Say it once more.
\end{tcolorbox}
\section[füllen]{füllen — to fill}

\label{lesson:GE-C146-fuellen}



\begin{quote}
Before the new one: say the German for to be missing, then the German for to celebrate.
\end{quote}

\subsection*{You'll want to know: füllen}
\textbf{füllen} — \textquotedblleft{}to fill\textquotedblright{}.

\begin{grammarlens}[title={What you've built}]
One more word for this chapter.
\end{grammarlens}

\subsection*{Your turn}
\begin{itemize}
\raggedright
  \item \emph{Say it:} \emph{füllen}
  \item \emph{Say it:} \emph{füllen}, once more
  \item \emph{Say it:} say \emph{füllen}
  \item \emph{Recall:} say the German for to be missing, then the German for to celebrate, then say \emph{füllen} again
\end{itemize}

\begin{tcolorbox}[breakable,colback=teal!4,colframe=teal!35!black,title={Before you move on}]
What does füllen mean? (\textquotedblleft{}To fill\textquotedblright{}.) Say it once more.
\end{tcolorbox}
\section[gefallen]{gefallen — to please, to be liked}

\label{lesson:GE-C146-gefallen}



\begin{quote}
Before the new one: say the German for to celebrate, then the German for to fill.
\end{quote}

\subsection*{You'll want to know: gefallen}
\textbf{gefallen} — \textquotedblleft{}to please, to be liked\textquotedblright{}.

\begin{grammarlens}[title={What you've built}]
One more word for this chapter.
\end{grammarlens}

\subsection*{Your turn}
\begin{itemize}
\raggedright
  \item \emph{Say it:} \emph{gefallen}
  \item \emph{Say it:} \emph{gefallen}, once more
  \item \emph{Say it:} say \emph{gefallen}
  \item \emph{Recall:} say the German for to celebrate, then the German for to fill, then say \emph{gefallen} again
\end{itemize}

\begin{tcolorbox}[breakable,colback=teal!4,colframe=teal!35!black,title={Before you move on}]
What does gefallen mean? (\textquotedblleft{}To please, to be liked\textquotedblright{}.) Say it once more.
\end{tcolorbox}
